% 06_data.tex --- Section 5: Data (empirical, post-fetch)
% ---------------------------------------------------------------------------
\section{Data}
\label{sec:data}

\subsection{Sources}
\label{sec:data-sources}

Aave V3 USDC reserve-rate history is pulled from the official Aave
protocol-subgraphs deployment on TheGraph's decentralized network,
subgraph ID
\texttt{Cd2gEDVeqnj\-Bn1hSeqFMitw8Q1\-iiyV9FYUZkLNRcL87g}
(entity \texttt{reserveParamsHistoryItems}, \texttt{liquidityRate} in
annualized RAY scaling, $10^{27}$). The Messari Compound V3 subgraph
indexes per-collateral markets rather than the base supply market;
USDC base-rate history is therefore reconstructed via batched
\texttt{eth\_call} on the Comet proxy
\texttt{0xc3d688B66703\-497DAA19211EEdff47f25384cdc3} at historical
block numbers, querying the
\texttt{getUtilization()} and \texttt{getSupplyRate(uint256)} view
functions. Gas prices come from Dune's \texttt{gas.gas\_price} table
at hourly median resolution; ETH/USD comes from CoinGecko's free
hourly close. Kink parameters per protocol per epoch are reconstructed
from on-chain \texttt{ReservesUpdated} (Aave) and
\texttt{Configurator} (Compound) events. The \texttt{fractal-defi}
substrate of \cite{krestenko2026basis} is used as the loader and
backtest harness.

\subsection{Window and Resampling}
\label{sec:window}

Eighteen months: 1~November~2024 through 30~April~2026, a
$13{,}104$-hour grid. The window spans two macro micro-regimes
(low-volatility H2~2024, normalizing 2025, and the present
higher-volatility 2026) without overlapping any protocol
re-deployment event. The Aave subgraph emits an event stream (median
$\sim 50$ events per hour, $\sim$$6.5\times 10^{5}$ raw rows over the
window) rather than uniform snapshots; the loader's \texttt{transform()}
resamples the stream to a right-closed hourly grid via
\texttt{df.resample("1H", label="right", closed="right").last()},
yielding \textbf{13{,}096 hourly Aave bars} ($99.9\%$ coverage). The
Compound RPC-reconstruction yields \textbf{4{,}900 hourly bars}
($37.4\%$ raw coverage of the same window); the gap is filled
incrementally as described in \cref{sec:coverage}.

\subsection{Empirical Spread Distribution}
\label{sec:empirical-spread}

On the $n = 4{,}894$ overlapping hours where both protocols have
observations, the cross-protocol supply-rate spread
$s(t) = \mathrm{APR}_{\text{Aave}}(t) - \mathrm{APR}_{\text{Compound}}(t)$
has the following marginal statistics: median
$-0.241$~pp, standard deviation $2.341$~pp, range
$[-13.32, +33.17]$~pp. The Aave-pays-more binary share is $43.4\%$;
Compound pays more on the typical hour by $0.24$~pp. The individual
APR distributions are likewise broad: Aave USDC supply APR ranges
$[1.54\%, 48.38\%]$ with median $3.82\%$; Compound USDC supply APR
ranges $[2.85\%, 23.65\%]$ with median $4.46\%$ (both quoted as
annualized continuous-compound rates).

Pooled marginals mask substantial structure. \Cref{tab:regime-quarters}
partitions the overlapping panel by calendar quarter:

\begin{table}[t]
\centering
\caption{Per-quarter cross-protocol spread structure on the overlapping
panel ($n = 4{,}894$ hours). Spread is
$\text{APR}_{\text{Aave}} - \text{APR}_{\text{Compound}}$ in
percentage points; ``Aave $>$ C'' is the share of hours on which Aave
pays more. The 2025 Q1$\to$Q2 jump from $25\%$ to $55\%$ is the
empirical regime shift motivating Branch~B of the forecaster.}
\label{tab:regime-quarters}
\begin{tabular}{lrrrr}
\toprule
Quarter & $n$ & spread median (pp) & spread std (pp) & Aave $>$ C \\
\midrule
2024 Q4 &  908 & $+0.10$ & $4.80$ & $51\%$ \\
2025 Q1 & 1336 & $-0.64$ & $1.43$ & $25\%$ \\
2025 Q2 & 1458 & $+0.22$ & $0.94$ & $55\%$ \\
2025 Q3 &  992 & $-0.07$ & $1.30$ & $45\%$ \\
2025 Q4 &  150 & $-0.15$ & $0.76$ & $39\%$ \\
2026 Q1 &   50 & $-0.23$ & $0.46$ & $38\%$ \\
\bottomrule
\end{tabular}
\end{table}

The direction-share jump from $25\%$ in 2025 Q1 to $55\%$ in 2025 Q2 is
a regime shift of the kind a dual-branch architecture with explicit
cross-protocol residual input is designed to detect; the four-fold
volatility contraction from $4.80$~pp (2024 Q4) to $0.94$~pp (2025 Q2)
also supplies natural anchors for the volatility-tertile ablation of
\cref{sec:regime}. Crucially, the binary direction split is close to
symmetric across the full panel ($43.4\%$ vs $56.6\%$), so a
forecast-driven edge cannot come from a structural bias toward either
protocol --- it must come from anticipating the crossovers, in keeping
with the cointegration finding of \cite{gudgeon2020defi} and against
the equilibrium-impossibility framing of \cite{halpern2024fairrates},
which concerns the existence of a single fair rate, not short-horizon
predictability of the realized spread.

\subsection{Data Coverage and Limitations}
\label{sec:coverage}

Raw coverage is asymmetric: Aave $99.9\%$ ($13{,}096 / 13{,}104$
hourly bars), Compound $37.4\%$ ($4{,}900 / 13{,}104$). The asymmetry
is structural: the Messari Compound V3 subgraph indexes per-collateral
markets and exposes \texttt{rates: None} on the base Comet market
entity, so hourly snapshots must be reconstructed via \texttt{eth\_call}
on archive RPCs --- and free-tier archive providers impose batch-size
caps near $100$ calls per request and per-day soft limits in the
thousands. The fetcher accordingly tops up the panel in an idempotent
append-mode across multiple free RPC providers (\texttt{publicnode}
and Ankr) until coverage saturates; the projected post-top-up coverage
on the joined panel is $\sim 70\%$. For sequence-length-168 forecaster
training, we restrict to dense subperiods in which both protocols have
recent observations, and forward-fill missing Compound rows up to a
$168$-hour ceiling on the joined panel, marking longer gaps as a
separate regime that is excluded from training and reported as edge
cases.

\subsection{Cleaning}
\label{sec:cleaning}

Outliers (rates outside $[0, 50\%]$, single-step relative jumps
$>5\times$) are flagged as oracle glitches and interpolated. The
sign-convention invariant
$\text{borrowing\_rate}\ge\text{lending\_rate}$ is enforced as a test
fixture (\texttt{tests/test\_sign\_convention.py}) at every timestamp
on both protocols, not as a trust statement about library version.
Z-score normalization uses training-split statistics only.

\subsection{Splits}
\label{sec:splits}

Strict block-number splits, to defend against subgraph-aggregation
leakage:
\begin{itemize}[leftmargin=*, itemsep=2pt]
\item \textbf{Train} (10 months): 2024-11-01 -- 2025-08-31,
  $\sim$$6{,}700$ hourly bars on both protocols.
\item \textbf{Validation} (4 months): 2025-09-01 -- 2025-12-31,
  $\sim$$2{,}920$ hourly bars on Aave but only $\sim$$290$
  Compound-dense hours pre-top-up; forecaster validation noise on the
  composite loss is expected to reflect this.
\item \textbf{Test} (4 months, held out): 2026-01-01 -- 2026-04-30,
  $\sim$$2{,}876$ hourly bars on Aave, $50$ Compound-real hours
  pre-top-up; the H1 evaluation on the full test window is
  Aave-driven in its initial form.
\end{itemize}
Real-data Compound sparsity means the H1 Sharpe-improvement claim on
the full $18$-month window is degraded relative to a fully-dense
panel; the evaluation protocol of \cref{sec:backtest} therefore
reports both the full-window result and a restricted result on dense
subperiods, and uses the Compound-rich 2024 Q4 sub-window
($n = 908$ overlapping hours, $51\%$ Aave-pays-more share) as an
additional H1 anchor.
